% Lösungen Einheit 2 von 2 – Einheit 2 (zusammenbau v0.1)
\einheitenkopf{Einheit 2 von 2 · Einheit 2}

\erg{11}{a) $1 - P(\overline{A} \cap \overline{B})$ – drei Felder, über das Gegenfeld. \quad b) $300 - 120 - (180 - 90) = 90$ \quad c) $P(W \cup T) = 58\,\% + 30\,\% - 21\,\% = 67\,\%$ (den Schnitt einmal abziehen); ebenso $1 - P(\overline{W} \cap \overline{T}) = 1 - 0{,}33$ \quad d) Erwachsene ohne Medaille: $30\,\% - 14\,\% = 16\,\%$ (die $14\,\%$ sind schon ein Anteil an allen); Kinder ohne Medaille: $0{,}5 \cdot 70\,\% = 35\,\%$; zusammen $16\,\% + 35\,\% = 51\,\%$ \quad e) „entweder $S$ oder $\overline{H}$“ heißt genau eines von beiden, also die Felder $S \cap H$ und $\overline{S} \cap \overline{H}$: $18\,\% + 33\,\% = 51\,\%$ (nicht das einschließende Oder, $78\,\%$); $\frac{27}{51} \approx 0{,}53$ – etwa die Hälfte, die Aussage stimmt.}
\erg{12}{a) Ida zieht den Schnitt nicht ab, er steckt in beiden Rändern. Richtig: $50\,\% + 40\,\% - 15\,\% = 75\,\%$. \quad b) Wer nicht „weder $A$ noch $B$“ ist, hat mindestens eines der Merkmale – genau das heißt „$A$ oder $B$“; zusammen füllen beide die ganze Tafel ohne Überschneidung.}
